\scanpage{29}

est finalement ramené aux $2$ questions suivantes

a) Déterminer $G=\mathrm{Isot}(S,S)$

b) Déterminer le sous-groupe \inserted{$H$} de $G$ formé des $h:S\longrightarrow S$ tels que $h\circ f|X\overset{\mathrm{isot}}{\simeq}f$.

On peut \unsure{évidemt|évidemment} supposer $S=S_\rho$ donc $G=G_\rho$, et $f=\varphi_\rho:X\longrightarrow S_\rho$, donc la question devient

a) Déterminer $G_{\unsure{S|\rho}}=\mathrm{Isot}(S_\rho,S_\rho)$ \inserted{$S_\rho$ surface de genre \deleted{\unsure{$g$}} $\unsure{\gamma}$ moins $\unsure{\nu|v}'$ pts \ldots}

b) Déterminer le sous-groupe $H_S\subset G_S$ des $h\in G_\rho$ tels que
\[
 \deleted{h|X\simeq}\qquad h\circ\varphi_\rho\simeq\varphi_\rho
\]
Tout dépend seulement du couple $X\subset S_\rho=S$, \deleted{\unsure{un}} surface orientée avec $\unsure{x|X}$ \inserted{connexe} un $X$ compact ($\in\mathrm{Ob}\,G_1$) plongé dedans.

La réponse à a) est classique (du moins si $S_\rho$ compact) \deleted{$G$ est} l'homomorphisme can
\[
 G\longrightarrow\mathrm{Aut\,ext}(\pi_1(S,s))
\]
est un Isom, et $G$ est donc le \uline{groupe de Teichmuller} de \deleted{\unsure{$S_\alpha$}} $S$. Reste à trouver (pour un $X\underset{\varphi}{\subset} S$ plongé dans $S$) quels sont les $h\in G$ tels que $h\circ\varphi\deleted{\unsure{x}}\underset{\mathrm{isot}}{\simeq}\varphi$.

\uline{Cas particulier} \inserted{$X$ connexe et} $\pi_1(X)\longrightarrow\pi_1(S)$ surjectif (sauf erreur, cela signifie \inserted{la partition \inserted{$(J_\lambda)$} de $J$ est discrète ($\Lambda=J$) et} que les $g_\lambda$ et les $\nu_\lambda$ sont nuls --- i.e. que $S_\rho$ est déduite de $\rho$ simplement en ``bouchant séparément chaque trou par une rondelle'' : plongements \uline{normaux}$^*$. Alors
% slanted
\margin{\unsure{pas} tout à fait vrai \unsure{sauf} peut-être. à vérifier\ldots.}
\[
 h\circ\varphi\underset{\mathrm{isot}}{\simeq}\varphi\quad\Longrightarrow\quad
 h\circ\varphi\underset{\mathrm{homot}}{\simeq}\varphi\qquad\text{donc}
\]
\diagram{\begin{tikzcd}[ampersand replacement=\&,row sep=small]
\pi_1(X)\arrow[r,"c"]\arrow[dr]\arrow[d,dash]\&\pi_1(S)\arrow[d,"\pi_1(h)"]\\
{}\&\pi_1(S)
\end{tikzcd}}
$\pi_1(h)\circ c=c$ mod autom \unsure{intérieurs|extérieurs}, donc $\pi_1(h)$ est un autom intérieur, $c$ étant surjectif, on \uline{gagne}.

Cela nous donne en particulier les plongts de \inserted{$X$ \uline{connexe}} dans $S^1$ (car tt \deleted{\unsure{plongt|plongement}} \inserted{\unsure{hom.|homéom.} orienté} de $S^1$ sur lui-mê est isotope à l'identité). Ils sont classifiés par les $\alpha$ tels que $\gamma(\alpha)$ [le genre de $\widetilde U_\alpha$ déduit de $U_\alpha$ par le procédé dit qui donne \deleted{\unsure{des}} un
% archive sheet number: 28
